\documentclass[11pt]{article}
\usepackage[margin=1in]{geometry}
\usepackage{amsmath,amssymb,booktabs,url,xcolor}
\usepackage[hidelinks]{hyperref}
\usepackage{graphicx}
\newcommand{\pending}[1]{\textcolor{red}{[PENDING: #1]}}

\title{Staged Upgrades of Two Learned Populations: Keeping a Learned Fleet's Benefit When the Signals It Relies On Are Upgraded\\[4pt]
\large DRAFT (05/10/2026)}
\date{}

\begin{document}
\maketitle

\begin{abstract}
Learned traffic signals and learned connected-vehicle fleets are deployed and upgraded on different schedules. We show, in SUMO, that a
learned eco-driving fleet that saves fuel under the old signal controller loses that benefit when the signal is upgraded to a learned
controller, and that the natural repairs fail: withdrawing the signal timing returns the fleet to its base driver, forecasting the learned
signal's timing costs the fleet's own vehicles time, and making the signal commit to its timing costs the whole network more than the fleet
recovers. The failure has a closed-loop cause that we isolate: once timing is published, eco-approaching vehicles stop being counted as
waiting by the halting-based learned signal, the signal holds the other direction's green longer, and the forecast's own target shifts.
We propose a \emph{two-sided interface} (TSI) that retrains neither population: the fleet receives a state-conditioned forecast that is
exact wherever the schedule is determined, and the frozen signal receives a demand adapter that counts slow-moving vehicles; both
interfaces are calibrated per deployed signal from training traffic only. In a pre-registered confirmation on five new signal seeds, four
new frozen fleets and fresh traffic, TSI reduces the fleet's fuel by $3.85\%$ $[3.26, 4.43]$, city time loss by $0.55$\,s/veh and human
time loss by $0.89$\,s/veh, with the fleet's time loss non-inferior ($-0.22$\,s/veh, one-sided $97.5\%$ upper bound $+0.04$ against a
$+0.5$ margin). On the same panel it matches a fleet retrained under the new signal on fuel and beats it on every time endpoint; both
interfaces are necessary, and the result holds for three receiver classes, fleet shares from 25 to 75\% (at 25\% the agent-time bound sits at the margin) and a 25\% demand increase. All
results are in simulation. On a second network with multi-green learned lights the mechanism replicates (the service alone is strongly
harmful and the adapter removes most of the harm), but TSI did not meet the time limits in two registered attempts.
\end{abstract}

\section{Introduction}
A connected-vehicle fleet trained to use signal phase and timing (SPaT) is useful only while the signals keep publishing timing it can act
on. Upgrading signals to learned adaptive controllers breaks this: the timing is no longer a static plan. Retraining the fleet each time a
city upgrades its signals, or retraining the signals each time a manufacturer updates its fleet, is often not possible; the two populations
are owned by different parties and released on different schedules. We ask whether a deployed learned fleet can keep its benefit when the
learned signals it relies on are upgraded, \emph{without retraining either population}.
\paragraph{Contributions.}
(i) A staged-upgrade benchmark with two learned populations (learned keep/switch signals; PPO, search and rule eco-driving fleets),
registered gates, and fresh-data confirmations (M1--M2).
(ii) A negative map of light-side and information-only repairs: commitments, gates, confidence and horizon filters, predictable splits and
per-decision commitment selection, each tested against the no-SPaT fallback; per-decision counterfactual credit is shown to be below the
simulator's chaotic noise floor.
(iii) A closed-loop diagnosis: publication of timing hides the fleet's demand from a halting-based learned signal, which lengthens greens and
shifts the forecast's own target (a performative feedback), while exact timing on a fixed schedule costs the fleet no time.
(iv) The two-sided interface TSI, confirmed on new signals, new fleets and fresh traffic, with component ablations, a same-panel comparison
with fleet retraining, and scope tests across receiver classes, penetration and demand.

\section{Setting}
SUMO grid4 (16 signals, 1\,800 veh/h, 50\% fleet share, membership fixed per vehicle id); lights FT / ACT / learned (RL-H: a DQN that
keeps or switches between two greens from halting and vehicle counts, holding with long phase durations); fleets: SUMO Krauss at
imperfection $\sigma$ (non-learned family), ACC, and the learned fleets of M1 (R9, a four-parameter search policy; R10, a PPO speed
residual on the lane-aware GLOSA rule). Determinism and provenance: PYTHONHASHSEED = 0, bit-identical always-SUMO across jobs (C-27).
\paragraph{Seeds and panels.} Learned lights: development seeds k0 / k1, confirmation seeds t4--t8 (never used to choose any M3 service).
Frozen fleets: development R10 s13 / s113, confirmation s213 / s313 / s413 / s513. Traffic: TRAIN 0--7 (fitting), DEV 600--611, SELECT
650--655, TUNE 660--663 (per-light calibration), confirmation 700--719 (first attempt) and 720--739 (second attempt), demand shift
740--749 (period 1.6\,s), scope 750--759; every panel was generated after a non-existence check, and no panel was used both to choose
and to confirm.

\section{Milestone M1: a learned fleet worth deploying}
\paragraph{Criterion (registered).} Agent fuel superiority vs always-SUMO with agent and human time-loss margins (+0.5 s/veh),
teleports and unfinished trips; strong baseline $\sigma^*$ chosen on SELECT; confirmation on four new policy seeds and fresh traffic.
\paragraph{Evidence so far.} No earlier recipe was confirmed (C-21 \dots C-31, corrected C-37); one development checkpoint (R5h\_t0)
passed on fresh traffic; the non-learned $\sigma$-family is a strong baseline ($\sigma=0.3$: $-2.5$ s/veh, $-2.2\%$ fuel).
\paragraph{Table 1 (fleet benefit; C-54, C-55).} Fresh traffic (20 seeds, generated after registration), FT, 50\% agents; learned
fleet = GLOSA-LD$\circ\sigma^*$, four policy parameters fitted by derivative-free search on training traffic (4 search seeds pooled).
Differences of the learned fleet, agent fuel / agent time loss / human time loss (95\% CI):
\begin{center}\small\begin{tabular}{lccc}\toprule
comparator & fuel (\%) & agent TL (s/veh) & human TL (s/veh)\\\midrule
always-SUMO & $-9.68$ [$-10.24$, $-9.00$] & $-3.48$ [$-3.77$, $-3.23$] & $-0.49$ [$-0.57$, $-0.41$]\\
$\sigma^*=0.3$ & $-7.58$ [$-8.14$, $-6.89$] & $-0.78$ [$-1.03$, $-0.58$] & $-0.36$ [$-0.45$, $-0.26$]\\
GLOSA (heuristic) & $-2.16$ [$-2.92$, $-1.35$] & $-3.03$ [$-3.88$, $-2.27$] & $-0.47$ [$-0.95$, $-0.07$]\\
GLOSA$\circ\sigma^*$ (no lane release) & $+0.16$ [$-0.57$, $+0.99$] & $-0.77$ [$-1.46$, $-0.16$] & $-0.30$ [$-0.65$, $+0.00$]\\
GLOSA-L$\circ\sigma^*$ (rule) & $-2.80$ [$-3.37$, $-2.06$] & $-0.18$ [$-0.41$, $+0.02$] & $-0.01$ [$-0.10$, $+0.08$]\\\bottomrule
\end{tabular}\end{center}
No teleports or unfinished trips for the learned fleet (GLOSA: 0.50, GLOSA$\circ\sigma^*$: 0.40 teleports per run). The fitted
parameters mainly learn where to release control (the last $\approx$15\,m before a needed lane change); the fuel gain over the rule is the
fuel the rule gives up by releasing early. Not RL; not compared with published eco-controllers; FT light only.
\paragraph{An RL fleet on the same panel (C-57).} A PPO speed residual on the rule (acting only where the rule commands; dense per-agent
cost; intervention-indexed discount, C-51), four new seeds pooled, against the rule: fuel $-1.70\%$ [$-2.33$, $-1.09$], agent time loss
$-0.04$ [$-0.29$, $+0.19$], human $+0.02$ [$-0.08$, $+0.12$]; against always-SUMO $-8.66\%$ / $-3.33$ / $-0.47$; against $\sigma^*$
$-6.54\%$ / $-0.64$ / $-0.33$; no teleports. It does not recover the fuel of GLOSA$\circ\sigma^*$ ($+1.29\%$ against it), which the fitted
release distance does.
\paragraph{Table 2 (SPaT ablation).} For the fitted GLOSA class, removing SPaT means the policy never commands and reduces to the
$\sigma^*$ fleet; the comparison against $\sigma^*$ above is therefore the class-level SPaT ablation ($-7.6\%$ fuel). It does not separate
information from representation. The PPO residual (R10, confirmed in M1, C-57) reduces to the $\sigma^*$ fleet exactly when no SPaT is
published (bit-identical rows, M2 check KM2-2), so its benefit against $\sigma^*$ is likewise the class-level SPaT ablation.
\paragraph{Diagnostics (appendix; C-39 \dots C-42).} D1: the ported phase reward (B) prefers physically better behaviour for large
differences but does not encode the time-loss constraints (it ranks an infeasible fuel-for-time policy far above always-SUMO); PPO from
scratch did not reach always-SUMO even on its training traffic. D2: a hand-written in-class controller (GLOSA target, no release) gives
$-6.7\%$ fuel at equal agent time but harms human cars; a behaviour clone of the heuristic GLOSA trades fuel for time in closed loop.
An evaluation bug (exit release missing in deterministic evaluation) was found by an identity check and fixed (C-41).
A residual on the in-class controller with a crawl-proof group cost (v9) failed its pilot gate; its confirmation was not evaluated (C-47).
Derivative-free policy search over a three-parameter GLOSA family on the deployment objective (fuel with time-loss penalties) returned
the untuned heuristic itself for the pilot search seed; the heuristic fails the $\sigma^*$ comparison on agent time loss (+1.8\,s/veh,
SELECT) and the human-time-loss margin on the width of its 5-seed bound (C-43).
\paragraph{An actuation pitfall shared by every commanded fleet (C-45).} With SUMO's \texttt{setSpeed}, every commanded fleet we
trained or wrote (GLOSA, an in-class controller, PPO, behaviour clones, earlier learned fleets) created teleports (2--17 per run), while
SUMO drivers at any $\sigma$ created none. Each inspected event is a lane-change deadlock: an agent car stopped at the end of a lane that does
not continue on its route. Releasing control whenever the car's current lane does not continue (a code predicate on SUMO's best-lanes
offset) removed these deadlocks on every panel tested (C-46). This is trace and intervention evidence in this SUMO version; without a
same-checkpoint factorial over the two interfaces we do not attribute the earlier learned fleets' time-loss failures to it in full (C-50).
\paragraph{A strong same-information rule (C-46).} GLOSA on the $\sigma^*$ base driver with lane-aware release (command only on a lane
that continues) is beneficial on fresh traffic (20 seeds): against $\sigma^*$, agent fuel $-5.0\%$ and agent / human time loss $-0.7$ /
$-0.4$\,s/veh; against always-SUMO, $-7.1\%$ and $-3.3$ / $-0.5$; no teleports. It is a fixed rule, not learned, not compared with
published eco-controllers, and without a SPaT ablation of its own it does not show that SPaT is necessary; it is the baseline a learned
fleet must beat on the same base and interface.
The learned four-parameter policy (R9) against this rule on fresh traffic is in Table~1 (C-54).

\section{Milestone M2: is there a compatibility problem?}
\paragraph{Design (target; partly run: M2a and M2b below; the adapted-fleet column is run as the retrained reference in M3; a learned $L_0$ has NOT been run).} Freeze the confirmed fleet $F_0$; lights $L_0$ (old) and $L_1$ (upgraded, learned), both beneficial in their
reference condition; matrix $L_0/F_0$, $L_1/F_0$, $L_0/F_1$, $L_1/F_1$ with absolute effects and difference-in-differences; SPaT
causal and correctly labelled (plan / guarantee / forecast with uncertainty).
\paragraph{M2a: SPaT service withdrawn or misused (C-59, C-60; fresh traffic, 20 seeds; four light seeds PAIRED with four fleet
seeds, not a factorial).} Scope of THIS panel: the old light is fixed-time (not a learned L0); it has no forecast arm (run later as M2b)
and no adapted-fleet column (run later as the retrained reference in M3). The shares of SPaT calls marked valid are calls, not measured forecast errors. The learned light is a large
upgrade over the fixed-time program with SUMO cars (city time loss $-33.7$\,s/veh, city fuel $-19.5\%$). Both confirmed fleets lose their
own benefit. If the learned light correctly publishes no plan, each fleet falls back to its base driver and becomes identical to
$\sigma^*$: the benefit drops from $-7.5\%$ (fitted fleet) and $-6.5\%$ (PPO fleet) agent fuel to exactly 0, with no harm. If the fleet
keeps deriving SPaT from the static program, which the learned light no longer follows (30--43\% of its SPaT calls are marked valid
outside the function's guaranteed scope; these are calls, not measured forecast errors), it is harmed against $\sigma^*$ under the same light: $+7.4\%$ / $+6.5\%$ fuel and $+1.7$ / $+1.6$\,s/veh agent time loss. In
absolute terms the agents still gain from the upgrade; what is lost is the fleet's benefit against its counterfactual drivers.
\paragraph{M2b: an empirical forecast from the learned light (C-61, C-71, C-72).} A conditional-quantile forecast of the learned light's
green durations does not return the benefit to the frozen fleet. Two facts explain why this is hard: the original forecast was 1\,s early in
the green branch on 100\% of points (a timing-convention error, fixed in a corrected provider; one yellow branch remains 1\,s early, C-84),
and the frozen fleet is sensitive to timing errors, with asymmetric fuel and time effects in an open-loop oracle-shift diagnostic (exact
timing $-6.6\%$ fuel; 2\,s early $-2.4\%$; 2\,s late $-4.2\%$ fuel but agent time loss $+1.5$\,s/veh; 8\,s late $+9.3\%$ and $+11.6$\,s/veh;
C-72). The diagnostic uses another fleet and schedule than the retrained fleet, so it warns of the trade-off but does not predict it.
\paragraph{Second network (C-81).} On an independent own-built network (SUMO \texttt{netgenerate} spider, 18 signals, six with three greens),
upgrading the static programs to a learned all-green controller cuts city time loss by 58.9\,s/veh. The frozen rule fleet (GLOSA-L$\circ\sigma^*$,
$-9.2\%$ agent fuel under the static programs) again loses its benefit: labelled $\to$ exactly $\sigma^*$ (no harm); naive $\to$ harmed
($+7.2\%$ fuel, agent TL $+2.7$, human TL $+0.6$, upper $+0.8$).
\begin{figure}[t]\centering\includegraphics[width=\linewidth]{figs_two_pop/fig1_m2_crossplay.pdf}
\caption{The frozen fleet's benefit against $\sigma^*$ under the same light, before (FT) and after the learned-light upgrade. Labelled bars are
exactly 0 (rows identical to $\sigma^*$). Crossed 95\% intervals; grid4 fresh 470--489 (four paired light/fleet seeds), spider fresh
1000--1019 (two light seeds). Script \texttt{scripts/CITY\_fig\_m2\_crossplay.py}.}\end{figure}
\paragraph{Known so far (not the target experiment).} Weak learned fleets regressed under light upgrades and a car-side reaction repair
helped (C-14, C-16); non-learned fleet upgrades ($\sigma$, ACC) were not harmed by three learned-light families (C-30, C-32 \dots C-34).
If $F_0$ shows no regression, the hazard hypothesis is rejected in this setting.

\section{Milestone M3: keeping the frozen fleet's benefit without retraining either population}
\paragraph{Endpoint (registered).} Agent fuel against R, the upgraded learned light publishing no SPaT (the frozen fleet then drives as
$\sigma^*$), with city, agent and human time-loss non-inferiority at $+0.5$\,s/veh (one-sided upper bounds), teleports and unfinished
trips; crossed bootstrap over light seed, fleet seed and traffic seed.

\subsection{What does not work (negative map)}
\paragraph{Light-side commitments.} Reward-shaped lights, fixed-window and ETA-triggered commitments, a predictable cycle and
receiver-triggered precision commitment (RPC) failed their DEV gates; on fresh traffic WC5 and RPC are worse than the fallback
($+2.98\%$ and $+9.26\%$ agent fuel, C-86). A $2\times2$ diagnostic of RPC (C-95) shows both suspected costs are real (the forced hold:
$-4.54\%$ fuel when removed; withdrawing the message at yellow onset: $-2.80\%$ when kept, $-2.66\%$ with the light schedule fixed), yet the best
cell is still $+3.06\%$: the commitment's own cost exceeds what the fleet recovers.
\paragraph{Choosing commitments decision by decision.} Re-simulating full episodes from identical prefixes with one intervention
(96 situations, five options, C-96) shows that one 5--15\,s commitment changes 640--900 of 1\,800 vehicles. Out of sample, no selection
rule (sensitivity, cost, both, or a learned regressor) beats a random option, and the best fixed option does not beat a permutation null
($p=0.29$; within-situation spread $0.235\%$ vs $0.040\%$ between options, C-98). Single-decision gains also do not add up: the
additive prediction from single edits is meaningless next to RPC's realized $+9.4\%$ (C-101). Candidates must therefore be evaluated as
whole policies over full episodes.
\paragraph{Information-only repairs.} Abstaining where the forecast leaves its support (0.11\% of the fleet's queries, C-97), confidence and
horizon gates (every gate moves along one fuel-for-agent-time frontier, C-104) and a predictable split controller tuned by CEM (it gives the
fleet $-7.9\%$ fuel from exact timing but is itself $+2.1$\,s/veh worse than the learned light, C-116) do not reach a feasible point.

\subsection{Diagnosis: publication hides demand from the learned light}
A state-conditioned direct time-to-green forecaster (an MLP quantile model on phase, age, queue and recent green history, fitted on training
traffic, C-99) lowers the query-weighted error by a third and gives the frozen fleet $-3.4\%$ fuel against R, but costs its vehicles
$+2.3$\,s/veh while human time falls ($-1.3$). Three diagnostics locate the cause.
(a) With SPaT published, mean halting on the served lanes falls from 2.9 to 1.6 and the learned light's greens lengthen (14.8 to 16.9\,s):
eco-approaching vehicles move above the $0.1$\,m/s halting threshold that the light observes and was trained on. Counting slow vehicles
(0.1--2\,m/s) as waiting in front of the frozen light cuts the agents' time cost by $0.8$\,s/veh (C-100).
(b) On the light's recorded schedule replayed open loop, exact timing gives the fleet $-5.6\%$ fuel \emph{and} $-0.37$\,s/veh; the time
cost is therefore a forecast error, not the fleet's own trade-off (C-107).
(c) Publishing the forecast shifts its own target: at the fleet's queries the error is $3.02$\,s when published vs $1.70$\,s when computed
but not used, and scored waits longer than 30\,s rise from $0.09\%$ to $3.34\%$ of queries (C-109); the demand adapter roughly halves this
sensitivity ($+0.73$ vs $+1.32$\,s, difference $-0.59$\,s, C-122). Across the five confirmation lights, the residual time cost grows with the residual green lengthening
(Figure~\ref{fig:mech}a, C-117).

\subsection{Method: a two-sided interface (TSI)}
Both learned populations keep their weights. (1) \emph{Fleet side}: a hybrid forecast, the program's exact answer where the schedule is
determined (the green after a yellow) and the state-conditioned forecaster elsewhere, served at the median. (2) \emph{Light side}: a demand
adapter in front of the frozen light's halting inputs, $q = h_{<0.1} + A\,n_{[0.1,2)} + B\,n_{[2,5)}$ on the served lanes. (3)
\emph{Calibration per deployed light from training traffic only}: the forecaster is fitted on eight training episodes of that light; $(A,B)$
is chosen on four tuning episodes by a pre-registered rule (agent, city and human time points $\le$ thresholds, then the most fuel). No
quantity of the evaluation traffic, lights or fleets enters calibration.

\subsection{Confirmation}
A first registered attempt with one global adapter chosen on two development lights (C-110) failed on agent time only: on five new lights,
four new fleets and fresh traffic (700--719), fuel $-4.56\%$, city $-0.20$, human $-1.23$, agent $+0.85$\,s/veh (upper $+1.27$, C-114). The
second and last registered attempt (TSI as above, per-light calibration; traffic 720--739; $\alpha=0.025$ one-sided because it is a second
attempt, Bonferroni bound $0.075$ over both) passed every limit (Table~\ref{tab:tsi}, Figure~\ref{fig:tsi}, C-118).
\begin{table}[t]\centering\small
\caption{TSI confirmation (C-118): five new learned-light seeds $\times$ four new frozen PPO fleets $\times$ 20 fresh traffic seeds, crossed
95\% intervals; the primary test uses the upper end (one-sided 97.5\%). F1-B: the fleet retrained under the new light (more update rights),
same panel.}\label{tab:tsi}
\begin{tabular}{lcccc}\toprule
vs R & agent fuel (\%) & city TL (s/veh) & agent TL (s/veh) & human TL (s/veh)\\\midrule
TSI & $-3.85$ [$-4.43$, $-3.26$] & $-0.55$ [$-0.67$, $-0.44$] & $-0.22$ [$-0.44$, $+0.04$] & $-0.89$ [$-1.10$, $-0.67$]\\
service only (HYB) & $-4.93$ & $+0.26$ & $+1.97$ & $-1.43$\\
adapter only (ADP) & $+0.29$ & $-0.17$ & $-0.09$ & $-0.25$\\
fleet retrained (F1-B) & $-3.18$ & $-0.17$ & $+0.28$ & $-0.62$\\\midrule
TSI $-$ F1-B & $-0.68$ [$-1.70$, $+0.11$] & $-0.38$ [$-0.59$, $-0.18$] & $-0.50$ [$-0.88$, $-0.09$] & $-0.27$ [$-0.44$, $-0.06$]\\\bottomrule
\end{tabular}\end{table}
\begin{figure}[t]\centering\includegraphics[width=\linewidth]{figs_two_pop/fig3_tsi_components.pdf}
\caption{TSI against the fallback R on the confirmation panel, with its two halves and the retrained fleet (95\% intervals; dashed: the
$+0.5$\,s/veh margin). Script \texttt{scripts/CITY\_fig\_tsi.py}.}\label{fig:tsi}\end{figure}
Both halves are necessary: the fleet-side service carries the fuel ($-4.12\%$ against the adapter alone) and the light-side adapter carries
the agents' time ($-2.19$\,s/veh against the service alone). The per-light chosen adapters are $(1.5,0.25)$, $(1,0.5)$, $(1,0.5)$,
$(1,0.25)$, $(1.5,0.5)$; agent time against R is within $[-0.47,+0.18]$\,s/veh on every light.
\begin{figure}[t]\centering\includegraphics[width=\linewidth]{figs_two_pop/fig4_tsi_mechanism.pdf}
\caption{Mechanism. (a) First TSI attempt (global adapter) and the service alone: per light, the agents' time cost grows with how much the
learned light lengthens its greens (C-117). (b) On a fixed schedule exact timing saves fuel and time; the forecast's errors cost time, and in
closed loop the light's reaction produces the fuel gain (C-100, C-107).}\label{fig:mech}\end{figure}

\subsection{Scope}
With the TSI design frozen (Figure~\ref{fig:scope}): a 25\% demand increase on fresh traffic, $-3.79\%$ fuel and $-0.27$\,s/veh agent
time (C-120); a rule fleet, $-1.84\%$ and $+0.03$; a search-fitted fleet, $-1.38\%$ and $+0.13$; fleet shares of 25 and 75\%, $-4.13\%$
/ $+0.07$ (upper $+0.50$, borderline) and $-4.75\%$ / $-0.58$ (C-121). Replacing the state forecaster by the age-only forecast with the same
adapter harms ($+6.72\%$; TSI is $9.87$ points better), so the state forecaster is necessary. The reverse upgrade order is safe in this
setting: a fleet retrained for the new light does no harm under the old light ($-2.23\%$ fuel, agent time $+0.15$, C-121).
\begin{figure}[t]\centering\includegraphics[width=\linewidth]{figs_two_pop/fig5_tsi_scope.pdf}
\caption{TSI with its design frozen across receiver classes, fleet shares and a demand shift (fresh traffic; 95\% intervals).}\label{fig:scope}\end{figure}

\section{Milestone M4: a second setting}
Spider network (own-built, \texttt{netgenerate}): a learned all-green light beats the static programs by 57\,s/veh (two seeds, C-76); the
rule fleet saves $-9.2\%$ fuel against $\sigma^*$ (C-77), with agent / human time-loss non-inferiority not established on fresh traffic
(upper $+0.55$ / $+0.52$ against the $+0.5$ margin, C-81);
the four-parameter search fleet (R9) does NOT beat the rule there (C-80); the M2 loss and harm replicate with the rule fleet (C-81); fixed-window commitment is again worse than the
no-SPaT fallback, with a different decomposition (no fleet contribution detected, C-82). The retrained-fleet check was not run there, because its registered condition (F1-A1 confirmed) failed. The
earlier RESCO cologne8 attempt failed at the light baseline and is kept as a negative (C-35, C-36).
\paragraph{TSI on the second network (C-123, C-124).} The same procedure (the forecaster generalized to two or three greens, the adapter
in front of the multi-green learned light's current / next-green inputs, per-light calibration on new tuning traffic) was confirmed twice on
fresh spider traffic with the rule fleet. The hidden-demand hazard is far larger there: the service alone raises the fleet's time loss by
$+59.5$\,s/veh and its fuel by $+20.7\%$. The adapter removes almost all of it (TSI $-$ service: $-58.1$\,s/veh), but a residual cost
remains: attempt 1, fuel $-1.66\%$ $[-1.99,-1.38]$, agent time loss $+1.33$\,s/veh (upper $+1.94$); attempt 2 with a wider adapter grid
($\alpha=0.025$), fuel $-0.31\%$, agent $+1.43$\,s/veh. TSI therefore does \emph{not} meet the M3 limits on the second network; its
mechanism replicates, and the two-parameter adapter appears too coarse for multi-green inputs (an untested explanation).

\section{Related work}
Joint or cooperative signal--vehicle learning retrains both populations (CoTV; HONEST-CAV, arXiv 2602.18740; SVCC). SPaT prediction for
actuated and adaptive signals is established (time-to-green learning; deep ensembles with confidence; HONEST-CAV's policy-aware
forecast); our forecaster is not new and matters only in combination. Observation adapters for frozen policies exist (DC-Ada; modular
world models for traffic-signal transfer, arXiv 2501.02548); ours is a two-parameter physical count transform calibrated against a
partner-induced shift. That published forecasts change their own target is the consistency problem of anticipatory route guidance (Bottom
and Ben-Akiva) and of performative prediction (Perdomo et al.; Drusvyatskiy and Xiao); we identify its traffic-signal instance and show that
a light-side input change reduces the sensitivity. The need for both interfaces matches a small-gain reading of the loop forecast $\to$ fleet
$\to$ light $\to$ target (Jiang, Teel and Praly; Dashkovskiy, R\"uffer and Wirth), and the adapter can be read as restoring an experiment
that Blackwell's ordering ranks as more informative for the light's decision (Blackwell 1953); both are motivation, not theorems about the
learned light.

\section{Limitations}
Simulation only (SUMO, HBEFA fuel). The method is confirmed on one network family (grid4); on the second network (spider, multi-green
lights, rule fleet) it reduced but did not remove the agents' time cost in two registered attempts (C-123, C-124). The confirmation that passed is a second attempt (family-wise bound $0.075$ with the first); the lights
and fleets of the two attempts are the same checkpoints, the traffic is fresh. The demand adapter needs per-lane speed sensing at the light
(speed bins on the served lanes). One learned-light family (keep/switch DQN on halting counts); a light trained on other inputs may not
hide demand the same way. The convention-corrected provider still has a 1\,s early yellow branch (C-84); the hybrid is exact only for the
green after a yellow (C-113). Per-decision analyses are limited by the simulator's chaotic divergence (C-98).

\section*{References (to be completed)}
\small
D.~Blackwell. Equivalent comparisons of experiments. \emph{Ann. Math. Statist.} 24(2):265--272, 1953.\\
J.~Bottom. Consistent anticipatory route guidance. PhD thesis, MIT, 2000.\\
S.~Dashkovskiy, B.~R\"uffer, F.~Wirth. An ISS small gain theorem for general networks. \emph{Math. Control Signals Systems} 19:93--122, 2007.\\
D.~Drusvyatskiy, L.~Xiao. Stochastic optimization with decision-dependent distributions. \emph{Math. Oper. Res.} 48(2):954--998, 2023.\\
Z.-P.~Jiang, A.~R.~Teel, L.~Praly. Small-gain theorem for ISS systems and applications. \emph{Math. Control Signals Systems} 7(2):95--120, 1994.\\
J.~Perdomo, T.~Zrnic, C.~Mendler-D\"unner, M.~Hardt. Performative prediction. \emph{ICML}, 2020.\\
HONEST-CAV, arXiv 2602.18740; Adaptive modular world models for traffic-signal control, arXiv 2501.02548; DHCP, arXiv 2509.03118;
time-to-green prediction (Genser et al., ETH); CoTV (Guo et al.); DC-Ada. \emph{Blackwell, Jiang--Teel--Praly and Drusvyatskiy--Xiao verified (05/10); the others still to be verified before submission.}
\normalsize

\appendix
\section{Validation and pitfalls (QA)}
Reward traces, sentinel readings, residual ratchet, horizon semantics, dominance labels (C-17, C-22, C-25, C-27, C-37, C-38).

\end{document}
